\chapter{Papers}
I can't believe how few papers I have written notes about, I should do better.
\section{Operator Splitting Methods: Numerical Solutions of Ordinary Differential Equations via Separation of Variables}

\subsection{General Idea}

It is deriving the mathematical underpinning of operator splitting from Koopman-Lie semigroup rather than physical systems. Generalizing it beyond Hamiltonian theory


\subsection{Notes On Introduction}
N/A

\subsection{Strongly Continuous and bi-continuous semigroups}


You know what, at the moment this paper may be beyond my graps.


\section{Splitting methods for differential equations}

\subsection{Notes on introduction and Abstract}

Some argue there are only ten big ideas in numerical analysis. They make a cool little connection to Descartes rules of logic and how splitting lays the foundation of science.

They then use a simple argument to derive strang splitting

Adjoint method refers to steping forward and then stepping backwardswith the inverse if its adjoint. Adjoint simply refers to that if $\psi_h = \varphi_1(\varphi_2)$ then $\psi^*_h = \varphi_2(\varphi_1)$

This is called time symmetric or self-adjoint method.

Some disadvantages of higher order splitting is that they go backwards in time which is troublesome in dissipative systems. extensively used in fields as distant as molecular dynamics, simulation of storage rings in particle accelerators, celestial mechanics, astronomy, quantum (statistical) mechanics, plasma physics, hydrodynamics, and Markov Chain Monte Carlo methods.

\subsection{Higher Order Splitting and Composition}




\section{Core Collapse Supernova Modeling: The Next Ten Years}


\section{Physical, numerical, and computational challenges of modeling neutrino transport in core-collapse supernovae}

\subsection{Structure Preserving}
